\section{特征值}

\begin{definition}
给定矩阵 $A \in M_n(\C)$,如果存在标量 $\lambda$ 和\textbf{非零}向量 $x$ 满足
\[
    Ax = \lambda x,
\]
那我们就称 $\lambda$ 是 $A$ 的一个特征值,$x$ 是这个特征值对应的一个特征向量.
\end{definition}

实际上,我们也常常把上面的式子写成 $Ax = x\lambda$ 的形式.当我们只写出一个特征值的时候可能还看不出什么,但当我们有足够多线性无关的特征向量的时候,我们就能写成
\[
    A[x_1,x_2,\cdots,x_n] = [x_1,x_2,\cdots,x_n]
    \begin{bmatrix}
        \lambda_1 & & & \\
        & \lambda_2 & & \\
        & & \ddots & \\
        & & & \lambda_n
    \end{bmatrix},
\]
进而就有
\[
    AQ = Q\Lambda \Rightarrow A = Q\Lambda Q^{-1}, \quad \Lambda = \operatorname{diag}(\lambda_1,\lambda_2,\cdots,\lambda_n).
\]
这也就是相似对角化的过程（如果 $A$ 可以相似对角化）.

如果刚好这个 $A$ 还是 Hermite 的,那么还会有一些特别的性质:比如它的特征向量能取为一组正交基.

\subsection{特征多项式}

由 $Ax = \lambda x$,我们可以得到:
\[
    (A - \lambda I)x = 0 \ \text{有非0解} \Rightarrow \det(A - \lambda I) = 0.
\]
于是我们定义
\[
    \begin{aligned}
        p(\lambda) &= \det(\lambda I - A) \\
        &= \lambda^n - \alpha_{n-1}\lambda^{n-1} + \cdots + (-1)^{n-1}\alpha_1\lambda + (-1)^n\alpha_0
    \end{aligned}
\]
为\textbf{特征多项式}.这里我们定义 $p(\lambda) = \det(\lambda I - A)$ 而不是 $\det(A - \lambda I)$ 的原因主要是保持 $\lambda^n$ 的系数为 1.

由代数基本定理,我们还能得到
\[
    p(\lambda) = \prod_{i=1}^n (\lambda - \lambda_i).
\]

\subsection{Vieta 定理}

回忆一下之前的式子:
\[
    \prod_{i=1}^n (\lambda - \lambda_i) = \lambda^n - \alpha_{n-1}\lambda^{n-1} + \cdots + (-1)^{n-1}\alpha_1\lambda + (-1)^n\alpha_0
\]
由韦达定理我们可以得到:
\[
    \begin{aligned}
        \alpha_{n-1} &= \lambda_1 + \lambda_2 + \cdots + \lambda_n, \\
        \alpha_{n-2} &= \lambda_1\lambda_2 + \lambda_1\lambda_3 + \cdots + \lambda_{n-1}\lambda_n, \\
        \vdotswithin{=} \\
        \alpha_0 &= \lambda_1\lambda_2\cdots\lambda_n.
    \end{aligned}
\]
那么如何求出 $\alpha_0,\alpha_1,\cdots,\alpha_{n-1}$ 的值呢?下面是一个比较重要的性质:

\begin{lemma}
$\alpha_k$ 等于 $A$ 的所有 $(n-k)$ 阶主子式之和.
\end{lemma}

% 原 Markdown 此处为图片证明,按要求未插入图片.

由此我们还可以推出两个比较常用的结论:
\begin{enumerate}
    \item $\alpha_0 = \lambda_1\lambda_2\cdots\lambda_n = \det(A)$；
    \item $\alpha_{n-1} = \lambda_1 + \lambda_2 + \cdots + \lambda_n = \operatorname{tr}(A)$.
\end{enumerate}

\subsection{代数重数和几何重数}

定义:
\[
    \begin{gathered}
        m_{\mathrm{alge}}(\lambda_i) = \sup\left\{k \in \mathbb{Z} \mid \lim_{\xi \to \lambda_i} p(\xi)(\xi - \lambda_i)^{-k} = 0 \right\}, \\
        m_{\mathrm{geom}}(\lambda_i) = \dim \operatorname{Ker}(\lambda_i I - A).
    \end{gathered}
\]

性质:$m_{\mathrm{alge}}(\lambda_i) \ge m_{\mathrm{geom}}(\lambda_i)$.

证明:

设 $A$ 对应 $\lambda_i$ 的几何重数为 $k$,即对应 $\lambda_i$ 最多有 $k$ 个线性无关的特征向量 $x_1,\cdots,x_k$.仿照上面的过程,我们有如下的式子:
\[
    \begin{gathered}
        Ax_1 = x_1\lambda_i \\
        Ax_2 = x_2\lambda_i \\
        \vdots \\
        Ax_k = x_k\lambda_i
    \end{gathered}
    \quad \Rightarrow \quad
    A[x_1,\cdots,x_k] = [x_1,\cdots,x_k]
    \begin{bmatrix}
        \lambda_i & & \\
        & \ddots & \\
        & & \lambda_i
    \end{bmatrix}
    = [x_1,\cdots,x_k]\lambda_i I_k
\]

我们可以把这 $k$ 个特征向量扩充为整个空间上的一组基 $\{x_1,\cdots,x_k,s_{k+1},\cdots,s_n\}$,并且记 $P = [x_1,\cdots,x_k,s_{k+1},\cdots,s_n]$.

那么我们有:
\[
    AP = P
    \begin{bmatrix}
        \lambda_i I_k & * \\
        0 & C
    \end{bmatrix}
\]
记
\[
    B = \begin{bmatrix}
        \lambda_i I_k & * \\
        0 & C
    \end{bmatrix},
\]
则
\[
    AP = PB \Rightarrow A = PBP^{-1}.
\]
由于这是一个相似变换,而相似变换不改变特征多项式,所以我们有
\[
    p_A(\lambda) = p_B(\lambda) = \det(\lambda I - B) = (\lambda - \lambda_i)^k p_C(\lambda).
\]
因此,矩阵 $A$ 与 $\lambda_i$ 对应的代数重数就是方程 $(\lambda - \lambda_i)^k p_C(\lambda)$ 在 $\lambda_i$ 处重根的个数.因为 $(\lambda - \lambda_i)^k$ 这个因子的存在,我们知道重根个数一定大于等于 $k$.因此 $m_{\mathrm{alge}} \ge k = m_{\mathrm{geom}}$.

\subsection{左特征向量和（左）特征值}

定义:
\[
    y^*A = \lambda y^*, \quad y \ne 0
\]
将上式左右分别取转置,我们有
\[
    A^*y = y\bar{\lambda}
\]
因此我们不难发现,左特征向量就是 $A^*$ 的特征向量,左特征值和 $A$ 的特征值基本没有区别.（好像差了一个共轭?但是 smy 没提,暂且存疑）

回忆上面相似对角化的过程,我们有
\[
    A[p_1,\cdots,p_n] = [p_1,\cdots,p_n]
    \begin{bmatrix}
        \lambda_1 & & \\
        & \ddots & \\
        & & \lambda_n
    \end{bmatrix}
\]
记
\[
    P = [p_1,\cdots,p_n], \quad \Lambda =
    \begin{bmatrix}
        \lambda_1 & & \\
        & \ddots & \\
        & & \lambda_n
    \end{bmatrix},
\]
则上式可以被写成
\[
    AP = P\Lambda
\]
两边同时左乘和右乘 $P^{-1}$,得
\[
    P^{-1}A = \Lambda P^{-1}.
\]
因此 $P^{-1}$ 也是由 $A$ 的左特征向量所构成的矩阵,即
\[
    P^{-1} =
    \begin{bmatrix}
        y_1^* \\
        \vdots \\
        y_n^*
    \end{bmatrix}.
\]
又由 $P^{-1}P = I$,我们有
\[
    y_i^*x_j = \delta_{ij}, \quad \text{其中} \delta_{ij} =
    \begin{cases}
        1, & i=j, \\
        0, & i \ne j.
    \end{cases}
\]
这就是左特征向量和右特征向量的\textbf{交叉正交性}.

\subsection{谱映射定理}

我们先思考这样一个问题:如果 $\lambda_i$ 是 $A$ 的特征值,那么 $A^2$ 的特征值会有什么性质呢?更进一步,$A^k$ 呢?

因为 $\lambda_i$ 是 $A$ 的特征值,设 $x_i$ 是这个特征值相对的特征向量,那么我们有
\[
    Ax_i = x_i\lambda_i \Rightarrow A^2x_i = A(Ax_i) = A(x_i\lambda_i) = x_i\lambda_i^2
\]
因此 $\lambda_i^2$ 也是 $A^2$ 的特征值.类似地,在 $k \ge 0$ 时,都有 $\lambda(A^k) = (\lambda(A))^k$.当 $\det(A) \ne 0$ 时,上式对 $k < 0$ 的情况也成立.

\textbf{谱映射定理:如果 $f$ 是解析函数,且 $\lambda$ 是 $A$ 的特征值,那么 $f(\lambda)$ 也是 $f(A)$ 的特征值.}

（这部分证明不难,所以我是让 AI 帮我写的）

证明:设 $A$ 的特征值为 $\lambda_1,\lambda_2,\dots,\lambda_n$,（这里按照代数重数重复计算,）我们要证明:$f(A)$ 的特征值为 $f(\lambda_1),f(\lambda_2),\dots,f(\lambda_n)$.

因为 $A \in \C^{n\times n}$,所以 $A$ 可以相似于一个上三角矩阵.也就是说,存在可逆矩阵 $P$,使得
\[
    A = PTP^{-1},
\]
其中
\[
    T =
    \begin{bmatrix}
        \lambda_1 & * & \cdots & * \\
        0 & \lambda_2 & \cdots & * \\
        \vdots & \vdots & \ddots & \vdots \\
        0 & 0 & \cdots & \lambda_n
    \end{bmatrix}.
\]
这里对角线上的元素就是 $A$ 的特征值,按照代数重数排列.

因为 $f$ 是解析函数,所以在合适区域内可以展开成幂级数:
\[
    f(z) = \sum_{m=0}^{\infty} a_m z^m.
\]
因此
\[
    f(A) = \sum_{m=0}^{\infty} a_m A^m.
\]
由于 $A = PTP^{-1}$,所以 $A^m = PT^mP^{-1}$.于是
\[
    f(A) = \sum_{m=0}^{\infty} a_m PT^mP^{-1} = P\left(\sum_{m=0}^{\infty} a_m T^m\right)P^{-1}.
\]
也就是说,
\[
    \boxed{f(A)=Pf(T)P^{-1}.}
\]
所以 $f(A)$ 与 $f(T)$ 相似,它们有相同的特征值.

现在只需要看 $f(T)$ 的特征值.

由于 $T$ 是上三角矩阵,所以 $T^m$ 也是上三角矩阵,并且 $T^m$ 的对角线元素是
\[
    \lambda_1^m,\lambda_2^m,\dots,\lambda_n^m.
\]
因此
\[
    f(T) = \sum_{m=0}^{\infty} a_m T^m
\]
仍然是上三角矩阵,而且它的第 $j$ 个对角线元素为
\[
    \sum_{m=0}^{\infty} a_m \lambda_j^m = f(\lambda_j).
\]
所以 $f(T)$ 是一个上三角矩阵,其对角线为
\[
    f(\lambda_1),f(\lambda_2),\dots,f(\lambda_n).
\]
而上三角矩阵的特征值就是其对角线元素.因此 $f(T)$ 的特征值为
\[
    f(\lambda_1),f(\lambda_2),\dots,f(\lambda_n).
\]
由于 $f(A)$ 与 $f(T)$ 相似,所以 $f(A)$ 也有这些特征值.因此
\[
    \boxed{\sigma(f(A)) = f(\sigma(A)).}
\]

\subsection{特征值与线性递推}

回忆 Fibonacci 数列的递推形式:
\[
    F_{n+2} = F_{n+1} + F_n
\]
我们可以把上面的递推式写成矩阵乘法的形式:
\[
    \begin{bmatrix}
        F_{n+2} & F_{n+1} \\
        F_{n+1} & F_n
    \end{bmatrix}
    =
    \begin{bmatrix}
        F_{n+1} & F_n \\
        F_n & F_{n-1}
    \end{bmatrix}
    \begin{bmatrix}
        1 & 1 \\
        1 & 0
    \end{bmatrix}.
\]
如果把这个矩阵乘法反复应用 $n$ 次,我们就能得到
\[
    \begin{bmatrix}
        F_{n+2} & F_{n+1} \\
        F_{n+1} & F_n
    \end{bmatrix}
    =
    \begin{bmatrix}
        F_2 & F_1 \\
        F_1 & F_0
    \end{bmatrix}
    \begin{bmatrix}
        1 & 1 \\
        1 & 0
    \end{bmatrix}^n
\]
其中,$\begin{bmatrix}1 & 1 \\ 1 & 0\end{bmatrix}^n$ 这一项我们只要知道它的特征值和特征向量,就能约化为 $(P\Lambda P^{-1})^n = P\Lambda^n P^{-1}$ 的形式.

对于 $\begin{bmatrix}1 & 1 \\ 1 & 0\end{bmatrix}$,它是方程 $\lambda^2 - \lambda - 1 = 0$ 所对应的 Frobenius 友矩阵,所以对每一个特征值 $\lambda$,$v(\lambda) = \begin{bmatrix}\lambda \\ 1\end{bmatrix}$ 都一定是它的特征向量.因此
\[
    \begin{bmatrix}
        1 & 1 \\
        1 & 0
    \end{bmatrix}^n
    =
    \begin{bmatrix}
        \lambda_1 & \lambda_2 \\
        1 & 1
    \end{bmatrix}
    \begin{bmatrix}
        \lambda_1^n & \\
        & \lambda_2^n
    \end{bmatrix}
    \begin{bmatrix}
        \lambda_1 & \lambda_2 \\
        1 & 1
    \end{bmatrix}^{-1}.
\]
利用 Frobenius 友型可以大大降低我们的计算量.

\subsection{Cayley-Hamilton 定理}

若 $p(\lambda) = \det(\lambda I - A)$,那么 $p(A) = 0$.
